\chapter{ماتریس و دترمینان}
\thispagestyle{empty}

\begin{quote}
	{\lotus
		ما [او و هالموس] فلسفه‌ی مشترکی درمورد جبر خطی داریم؛ فارغ از پایه فکر می‌کنیم، فارغ از پایه می‌نویسیم، اما هنگامی که کار به جای باریک می‌رسد درِ اتاق کارمان را می‌بندیم و با عصبانیت آن را با ماتریس‌ها محاسبه می‌کنیم. \quad - کاپلانسکی%
		\LTRfootnote{ Irving Kaplansky}
	}
\end{quote}

\section{ماتریس}
 هر آرایش مستطیلی از اعداد، شامل تعدادی سطر و ستون یک ماتریس نامیده می‌شود. هر عدد واقع در هر ماتریس را درایه‌ی آن ماتریس می‌نامیم. معمولاً ماتریس‌ها را با حروف بزرگ مانند $A$، $B$، $C$ و ... نام‌گذاری می‌کنیم.

در حالت کلی، اگر ماتریسی چون $A$ دارای $m$ سطر و $n$ ستون باشد، می‌نویسیم
$A_{m \times n}$
و می‌خوانیم $A$ ماتریسی از مرتبه‌ی $m$ در $n$ است. برای هر درایه‌ی ماتریس دو اندیس در نظر می‌گیریم که اندیس سمت چپ جای سطر و اندیس سمت راست جای ستون آن درایه را مشخص می‌کند،
$a_{ij}$
یعنی درایه‌ی روی سطر $i$ و ستون $j$.
\[
A = \begin{bmatrix}
	a_{11} & a_{12} & \cdots & a_{1n} \\
	a_{21} & a_{22} & \cdots & a_{2n} \\
	\vdots & \vdots & \ddots & \vdots \\
	a_{m1} & a_{m2} & \cdots & a_{mn}
\end{bmatrix}
\]

\begin{example}
	ماتریس $B$ ماتریسی شامل دو سطر و سه ستون است. این ماتریس دارای
	$2 \times 3 = 6$
	درایه است و مثلاً عدد
	$\sqrt{3}$
	درایه‌ی روی سطر اول و ستون اول است و درایه‌ی
	$(-4)$
	روی سطر دوم و ستون سوم قرار دارد.
	\[
	B = \begin{bmatrix}
		\sqrt{3} & 0 & \dfrac{2}{5} \\
		\pi & 6 & -4
	\end{bmatrix}
	\]
\end{example}

\begin{example}
	ماتریس
	$A=[a_{ij}]_{2 \times 2}$
	را در نظر بگیرید. اگر برای
	$i=j$
	داشته باشیم
	$a_{ij}=9$
	و برای
	$i>j$
	داشته باشیم
	$a_{ij}=-4$
	 و برای
	 $i<j$
	 داشته باشیم
	 $a_{ij}=6$،
	 در این صورت ماتریس $A$ را با درایه‌هایش نمایش دهید.
	 \\
	 پاسخ:
	 \[
	 A = \begin{bmatrix}
	 	9 & 6 \\
	 	-4 & 9
	 \end{bmatrix}
	 \]
\end{example}
\bigskip
دو ماتریس هم‌مرتبه‌ی
$A=[a_{ij}]_{m \times n}$
و
$B=[b_{ij}]_{m \times n}$
را مساوی گوییم هرگاه درایه‌های نظیر به نظیر آن‌ها با هم برابر باشند.
$$\forall i,j : a_{ij} = b_{ij} \Leftrightarrow [a_{ij}] = [b_{ij}]$$

\begin{example}
	اگر دو ماتریس
	$A = \begin{bmatrix}
		x-y & 8 \\
		4 & z-1
	\end{bmatrix}$
	و
	$B = \begin{bmatrix}
		2 & x+y \\
		4 & 6
	\end{bmatrix}$
	مساوی باشند،
	$(x+y+z)$
	را بیابید. \\
	\textbf{پاسخ:}
	\[
	A=B \Rightarrow
	\begin{cases}
		x-y=2 \\
		x+y=8 \\
		z-1=6
	\end{cases}
	\Rightarrow
	x=5 , y=3 , z=7
	\Rightarrow
	x+y+z=15
	\]
\end{example}

\subsection{ماتریس‌های خاص}
\textbf{ماتریس مربعی:}
اگر در ماتریس $A$، تعداد سطرها با تعداد ستون‌ها برابر و مساوی $n$ باشد، $A$ را یک ماتریس مربعی از مرتبه‌ی $n$ می‌نامیم. اگر
$i=j$
در این صورت درایه‌ی
$a_{ij}$
روی
\textbf{قطر اصلی}
ماتریس قرار دارد.
\begin{example}
	\[
	A = \begin{bmatrix}
		5 \\
	\end{bmatrix}_{1 \times 1}
	= 5
	\qquad
	B = \begin{bmatrix}
		1 & 2 \\
		-1 & 0
	\end{bmatrix}_{2 \times 2}
	\qquad
	C = \begin{bmatrix}
		2 & 1 & 3 \\
		-4 & 5 & 0 \\
		3 & 7 & 8
	\end{bmatrix}_{3 \times 3}
	\]
\end{example}
\smallskip
\textbf{ماتریس سطری:}
اگر ماتریس $A$ فقط دارای یک سطر باشد، آن را یک ماتریس سطری می‌نامیم.
\begin{example}
	\[
	A = \begin{bmatrix}
		11
	\end{bmatrix}_{1 \times 1}
	=11
	\qquad
	B = \begin{bmatrix}
		3 & 7
	\end{bmatrix}_{1 \times 2}
	\qquad
	C = \begin{bmatrix}
		2 & -1 & 4 & 6
	\end{bmatrix}_{1 \times 4}
	\]
\end{example}
\smallskip
\textbf{ماتریس ستونی:}
اگر ماتریس $A$ فقط دارای یک ستون باشد، آن را یک ماتریس ستونی می‌نامیم.
\begin{example}
	\[
	A = \begin{bmatrix}
		23
	\end{bmatrix}_{1 \times 1}
	=23
	\qquad
	B = \begin{bmatrix}
		0 \\
		\sqrt{60} \\
		-53 \\
		\pi
	\end{bmatrix}_{4 \times 1}
	\]
\end{example}
\smallskip
\textbf{ماتریس قطری:}
ماتریس قطری، ماتریسی است مربعی که تمام درایه‌های غیرواقع بر قطر اصلی آن صفر باشند. درایه‌های واقع بر قطر اصلی می‌توانند صفر باشند یا نباشند.
\begin{example}
	\[
	A = \begin{bmatrix}
		2 & 0 \\
		0 & 7
	\end{bmatrix}
	\qquad
	B = \begin{bmatrix}
		-1 & 0 & 0 \\
		0 & 0 & 0 \\
		0 & 0 & 8
	\end{bmatrix}
	\qquad
	C = \begin{bmatrix}
		0 & 0 & 0 \\
		0 & 0 & 0 \\
		0 & 0 & 0
	\end{bmatrix}
	\]
\end{example}
\smallskip
\textbf{ماتریس اسکالر:}
اگر ماتریسی قطری باشد و تمام درایه‌های روی قطر اصلی آن با هم برابر باشند، آن را یک ماتریس اسکالر می‌نامیم.
\begin{example} \label{exmatsca}
	\[
	A = \begin{bmatrix}
		3
	\end{bmatrix}
	\qquad
	B = \begin{bmatrix}
		1 & 0 \\
		0 & 1
	\end{bmatrix}
	\qquad
	C = \begin{bmatrix}
		5 & 0 & 0 \\
		0 & 5 & 0 \\
		0 & 0 & 5
	\end{bmatrix}
	\]
\end{example}
\smallskip
\textbf{ماتریس واحد یا همانی:}
ماتریس اسکالری که درایه‌های قطر اصلی آن یک باشد، ماتریس واحد یا همانی نامیده می‌شود. ماتریس همانی مرتبه‌ی $n$ را با نماد
$I_n$
نشان می‌دهند. در مثال
\ref{exmatsca}،
ماتریس $B$ یک ماتریس واحد یا همانی از مرتبه‌ی ۲ می‌باشد.
\\
\textbf{ماتریس صفر:}
ماتریس صفر، ماتریسی است که همه‌ی درایه‌های آن صفر باشند. ماتریس صفر را با نماد
$O$
یا
$\overline{O}$
نشان می‌دهیم.
\begin{example}
	$$O_{1 \times 2} =
	\begin{bmatrix}
		0 & 0
	\end{bmatrix}$$
\end{example}
\smallskip
\textbf{ماتریس پایین‌مثلثی:}
اگر درایه‌های بالای قطر اصلی ماتریس $A$ برابر صفر باشند، $A$ را ماتریس پایین‌مثلثی می‌نامیم. اگر درایه‌های قطر اصلی نیز برابر صفر باشند، آن ماتریس را پایین‌مثلثی اکید می‌نامیم.
\begin{example}
	ماتریس
	$
	A = \begin{bmatrix}
		1 & 0 & 0 \\
		4 & 2 & 0 \\
		5 & 3 & 7
	\end{bmatrix}
	$
	یک ماتریس پایین‌مثلثی از مرتبه‌ی ۳ است.
\end{example}
\smallskip
\textbf{ماتریس بالامثلثی:}
اگر درایه‌های پایین قطر اصلی ماتریس $A$ برابر صفر باشند، $A$ را ماتریس بالامثلثی می‌نامیم. اگر درایه‌های قطر اصلی نیز برابر صفر باشند، آن ماتریس را بالامثلثی اکید می‌نامیم.
\begin{example}
	ماتریس
	$
	B = \begin{bmatrix}
		1 & 9 & 2 \\
		0 & 3 & 5 \\
		0 & 0 & 6
	\end{bmatrix}
	$
	یک ماتریس بالامثلثی از مرتبه‌ی ۳ است.
\end{example}
\smallskip
\textbf{ترانهاده‌ی یک ماتریس:}
فرض کنیم
$A=[a_{ij}]_{m \times n}$
یک ماتریس باشد. ترانهاده‌ی $A$ که آن را با $A^t$ یا $A'$ نمایش می‌دهیم، برابر است با
$A^t=[b_{ij}]_{n \times m}$
و به صورت زیر تعریف می‌شود:
$$b_{ij} = a_{ji}$$

\begin{example}
	اگر
	$
	A = \begin{bmatrix}
		2 & \sqrt{2} \\
		-1 & 3 \\
		0 & 1
	\end{bmatrix}
	$،
	آنگاه
	$
	A^t = \begin{bmatrix}
		2 & -1 & 0 \\
		\sqrt{2} & 3 & 1
	\end{bmatrix}
	$.
\end{example}
\smallskip
\textbf{ماتریس متقارن:}
فرض کنیم
$A=[a_{ij}]_{n \times n}$
یک ماتریس باشد. $A$ را متقارن گوییم هرگاه
$A^t=A$.

\begin{example}
	مقادیر $x$ و $y$ و $z$ را چنان تعیین کنید که ماتریس $A$ متقارن باشد.
	\[
	A = \begin{bmatrix}
		1 & 3 & -2 \\
		2x-1 & 3 & y-1 \\
		3y+1 & -x & -4
	\end{bmatrix}
	\]
	\textbf{پاسخ:}
	\[
	\begin{cases}
	2x-1 = 3 \\
	3y+1 = -2 \\
	y-1 = -x
	\end{cases}
	\Rightarrow
	x=2 ,\: y=-1
	\]
	
\end{example}
\smallskip
\textbf{ماتریس پادمتقارن:}
فرض کنیم
$A=[a_{ij}]_{n \times n}$
یک ماتریس باشد. $A$ را پادمتقارن گوییم هرگاه
$A^t=-A$.

\begin{example}
	$A$
	یک ماتریس پادمتقارن مرتبه‌ی $n$ است. حاصل‌جمع همه‌ی درایه‌های $A$ چند است؟ \\
	\textbf{پاسخ:}
	درایه‌‌های واقع بر قطر اصلی همگی صفر هستند و درایه‌های غیر قطر اصلی نیز دوبه‌دو قرینه‌اند، پس مجموع همه‌ی درایه‌ها صفر است.
\end{example}

\subsection{جمع و ضرب در ماتریس‌ها}
\subsubsection*{جمع ماتریس‌ها}
برای جمع یا تفاضل دو ماتریس هم‌مرتبه‌ی $A$ و $B$ کافی است درایه‌های نظیر به نظیر دو ماتریس را با هم جمع یا از هم کم کرد.
$$A=[a_{ij}]_{m \times n} ,\ B=[b_{ij}]_{m \times n} \Rightarrow A \pm B = [a_{ij}] \pm [b_{ij}] = [a_{ij} \pm b_{ij}]$$

\begin{example}
	دو ماتریس
	$
	A = \begin{bmatrix}
		7 & 3 & 5 \\
		-5 & 2 & 0
	\end{bmatrix}
	$
	و
	$
	B = \begin{bmatrix}
		3 & 0 \\
		1 & 6 \\
		4 & 3
	\end{bmatrix}
	$
	را نمی‌توان با هم جمع کرد، چون هم‌مرتبه نیستند. ولیکن مجموع دو ماتریس
	$
	C = \begin{bmatrix}
		3 & 1 & 5 \\
		0 & 6 & 8
	\end{bmatrix}
	$
	و
	$
	D = \begin{bmatrix}
		4 & 5 & 6 \\
		9 & 2 & 1
	\end{bmatrix}
	$
	برابر است با ماتریس
	$
	C+D = \begin{bmatrix}
		7 & 6 & 11 \\
		9 & 8 & 9
	\end{bmatrix}
	$.
\end{example}

\begin{example}
	اگر
	$A=\begin{bmatrix}
		1 & 2 \\
		-3 & 5
	\end{bmatrix}$
	و
	$B=\begin{bmatrix}
		2 & -1 \\
		3 & 4
	\end{bmatrix}$،
	آنگاه
	$(A+B)^t$
	و
	$A^t + B^t$
	را به دست آورید. \\
	\textbf{پاسخ:}
	\[
	A+B =
	\begin{bmatrix}
		1 & 2 \\
		-3 & 5
	\end{bmatrix}
	+
	\begin{bmatrix}
		2 & -1 \\
		3 & 4
	\end{bmatrix}
	=
	\begin{bmatrix}
		3 & 1 \\
		0 & 9
	\end{bmatrix}
	\]
	
	\[
	(A+B)^t =
	\begin{bmatrix}
		3 & 0 \\
		1 & 9
	\end{bmatrix}
	\]
	
	\[
	A^t + B^t =
	\begin{bmatrix}
		1 & -3 \\
		2 & 5
	\end{bmatrix}
	+
	\begin{bmatrix}
		2 & 3 \\
		-1 & 4
	\end{bmatrix}
	=
	\begin{bmatrix}
		3 & 0 \\
		1 & 9
	\end{bmatrix}
	\]
\end{example}

\subsubsection*{ضرب یک عدد در یک ماتریس}
برای ضرب یک عدد در ماتریسی چون $A$ آن عدد را در تمام درایه‌های ماتریس ضرب می‌کنیم.
$$A=[a_{ij}]_{m \times n} ,\ r \in \mathbb{R} \Rightarrow rA = r[a_{ij}]_{m \times n} = [ra_{ij}]_{m \times n}$$

\begin{example}
	\[
	-2 \times
	\begin{bmatrix}
	2 & 6 & -4 \\
	3 & -2 & 0
	\end{bmatrix}
	=
	\begin{bmatrix}
		-4 & -12 & 8 \\
		-6 & 4 & 0
	\end{bmatrix}
	\]
\end{example}
\bigskip
\textbf{قرینه‌ی ماتریس:}
اگر
$A=[a_{ij}]_{m \times n}$
ماتریسی دلخواه باشد، قرینه‌ی ماتریس $A$ را با
$(-A)$
نمایش داده و از ضرب
$(-1)$
در ماتریس $A$ به دست می‌آید.

\subsubsection*{خواص مهم جمع ماتریس‌ها و ضرب عدد در ماتریس}
اگر $A$، $B$ و $C$ ماتریس‌هایی هم‌مرتبه و $r$ و $s$ اعدادی حقیقی باشند خواص زیر قابل اثبات‌اند:
\LR{
\begin{enumerate}
	\item $A+B = B+A$
	\item $A+(B+C) = (A+B)+C$
	\item $A+O = O+A = A$
	\item $A+(-A) = (-A)+A = O$
	\item $r(A \pm B) = rA \pm rB$
	\item $(r \pm s)A = rA \pm sA$
	\item $rA=rB ,\ r \neq 0 \Rightarrow A=B$
	\item $A=B \Rightarrow rA=rB$
\end{enumerate}}
\vspace*{-1cm}

\subsubsection*{ضرب ماتریس در ماتریس}
فرض کنیم
$A=[a_{ij}]_{m \times n}$
و
$B=[b_{ij}]_{n \times p}$.
حاصل‌ضرب $A$ در $B$ یعنی
$AB=[c_{ij}]_{m \times p}$
و به صورت زیر تعریف می‌شود:
$$c_{ij} = \sum_{k=1}^{n}a_{ik}b_{kj}$$
توجه کنید که تعداد ستون‌های $A$ برابر تعداد سطرهای $B$ است.

\begin{example}
	اگر
	$A=
	\begin{bmatrix}
		1 & 2 \\
		3 & 4 \\	
	\end{bmatrix}
	$
	و
	$B=
	\begin{bmatrix}
		5 & 6 \\
		7 & 8 \\	
	\end{bmatrix}
	$،
	حاصل $AB$ و $BA$ را حساب کنید. \\
	\textbf{پاسخ:}
	\[
	AB=
	\begin{bmatrix}
		1 & 2 \\
		3 & 4 \\	
	\end{bmatrix}
	\begin{bmatrix}
		5 & 6 \\
		7 & 8 \\	
	\end{bmatrix}
	=
	\begin{bmatrix}
		(1 \times 5)+(2 \times 7) & (1 \times 6)+(2 \times 8) \\
		(3 \times 5)+(4 \times 7) & (3 \times 6)+(4 \times 8) \\	
	\end{bmatrix}
	=
	\begin{bmatrix}
		19 & 22 \\
		43 & 50 \\	
	\end{bmatrix}
	\]
	\[
	BA=
	\begin{bmatrix}
		5 & 6 \\
		7 & 8 \\	
	\end{bmatrix}
	\begin{bmatrix}
		1 & 2 \\
		3 & 4 \\	
	\end{bmatrix}
	=
	\begin{bmatrix}
		(5 \times 1)+(6 \times 3) & (5 \times 2)+(6 \times 4) \\
		(7 \times 1)+(8 \times 3) & (7 \times 2)+(8 \times 4) \\	
	\end{bmatrix}
	=
	\begin{bmatrix}
		23 & 34 \\
		31 & 46 \\	
	\end{bmatrix}
	\]
\end{example}

\begin{example}
	مثال نقضی برای این گزاره بیابید: اگر حاصل‌ضرب دو ماتریس برابر صفر شود، حداقل یکی از آن‌ها صفر است. \\
	\textbf{پاسخ:}
	\[
	A=
	\begin{bmatrix}
		2 & -1 \\
		4 & -2
	\end{bmatrix},
	\qquad
	B=
	\begin{bmatrix}
		1 & 5 \\
		2 & 10
	\end{bmatrix}
	\]
	\[
	AB=
	\begin{bmatrix}
		2-2 & 10-10 \\
		4-4 & 20-20
	\end{bmatrix}
	=
	\begin{bmatrix}
		0 & 0 \\
		0 & 0
	\end{bmatrix}
	= O_2
	\]
\end{example}

\begin{example}
	با یک مثال نقض نشان دهید که در حالت کلی از تساوی
	$AB=AC$
	نمی‌توان نتیجه گرفت
	$B=C$.
	\\ \textbf{پاسخ:}
\end{example}

\begin{example}
	اگر
	$A=
	\begin{bmatrix}
		1 & -1 \\
		0 & 2
	\end{bmatrix}$،
	تمام ماتریس‌های
	$2 \times 2$
	مانند $B$ را بیابید که داشته باشیم
	$AB=BA$. \\
	\textbf{پاسخ:}
	فرض کنیم
	$B=
	\begin{bmatrix}
		a & b \\
		c & d
	\end{bmatrix}$.
	\[
	AB =
	\begin{bmatrix}
		1 & -1 \\
		0 & 2
	\end{bmatrix}
	\begin{bmatrix}
		a & b \\
		c & d
	\end{bmatrix}
	=
	\begin{bmatrix}
		a-c & b-d \\
		2c & 2d
	\end{bmatrix}
	\]
	\[
	BA =
	\begin{bmatrix}
		a & b \\
		c & d
	\end{bmatrix}
	\begin{bmatrix}
		1 & -1 \\
		0 & 2
	\end{bmatrix}
	=
	\begin{bmatrix}
		a & -a+2b \\
		c & -c+2d
	\end{bmatrix}
	\]
	\[
	AB=BA \Rightarrow
	\begin{cases}
		a-c=a \\
		2c=c \\
		b-d=-a+2b \\
		2d=-c+2d
	\end{cases}
	\Rightarrow
	c=0 , \ b=a-d
	\]
	بنابراین $a$ و $d$ اعدادی دلخواه هستند و تمام ماتریس‌های به شکل زیر پاسخ مسئله هستند.
	\[
	B=
	\begin{bmatrix}
		a & a-d \\
		0 & d
	\end{bmatrix}
	\]
\end{example}

\subsubsection*{خواص مهم ضرب ماتریس‌ها}
فرض کنیم
$A=[a_{ij}]_{m \times n}$،
$B=[b_{ij}]_{n \times p}$،
$C=[c_{ij}]_{p \times q}$ و
$D=[d_{ij}]_{n \times n}$.
در این صورت:
\LR{
	\begin{enumerate}
		\item $A(BC)=(AB)C$
		\item $A(B+C)=AB+BC$
		\item $(A+B)C=AC+BC$
		\item $DI_n=I_nD=D$
		\item $DO_n=O_nD=O_n$
\end{enumerate}}

فرض کنیم
$A=[a_{ij}]_{n \times n}$
و
$B=[b_{ij}]_{n \times n}$
دو ماتریس باشند و $r$ یک عدد حقیقی. در این صورت:
\LR{\begin{enumerate}
		\item $(A+B)^t = A^t + B^t$
		\item $(rA)^t = rA^t$
		\item $(AB)^t = B^t A^t$
		\item $(A^t)^t = A$
\end{enumerate}}

فرض کنیم
$A=[a_{ij}]_{n \times n}$.
در این صورت:
\[
\left[ \frac{1}{2}(A+A^t) \right]^t = \frac{1}{2}(A+A^t)^t = \frac{1}{2}(A^t+(A^t)^t) = \frac{1}{2}(A^t+A) = \frac{1}{2}(A+A^t)
\]
\[
\left[ \frac{1}{2}(A-A^t) \right]^t = \frac{1}{2}(A-A^t)^t = \frac{1}{2}(A^t-(A^t)^t) = \frac{1}{2}(A^t-A) = -\frac{1}{2}(A-A^t)
\]
در نتیجه
$\frac{1}{2}(A+A^t)$
ماتریس متقارن و
$\frac{1}{2}(A-A^t)$
ماتریس پادمتقارن است. اکنون با توجه به
$$A=\frac{1}{2}(A+A^t)+\frac{1}{2}(A-A^t)$$
نتیجه می‌گیریم که هر ماتریس مربعی را می‌توان به صورت مجموع یک ماترس متقارن و یک ماتریس پادمتقارن نوشت.

\textbf{تعویض‌پذیری:}
اگر $AB$ و $BA$ با هم برابر باشند، می‌گوییم ماتریس‌های $A$ و $B$ تعویض‌پذیر هستند. اگر $A$ و $B$ دو ماتریس مربعی هم‌مرتبه و تعویض‌پذیر باشند، آنگاه تمام اتحادهای جبری را درمورد آن دو ماتریس می‌توان به کار برد. همچنین اگر $m$ و $n$ دو عدد طبیعی باشند، آنگاه
$A^mB^n = B^nA^m$،
یعنی هر توانی از $A$ با هر توانی از $B$ تعویض‌پذیر است. اگر دو ماتریس تعویض‌پذیر نباشند، در نوشتن عبارت‌های جبری باید ترتیب ضرب حفظ شود. همچنین برای
$a,b,n,m \in \mathbb{N}$
داریم:
$$(A^a)^n \times (A^b)^m = (A^b)^m \times (A^a)^n= A^{an+bm}$$

\begin{example}
	اگر
	$A^3=
	\begin{bmatrix}
		3 & 2 \\
		2 & 1
	\end{bmatrix}$ و
	$A^4=
	\begin{bmatrix}
		5 & 3 \\
		3 & 2
	\end{bmatrix}$
	باشد، ماتریس
	$A^{11}$
	را به دست آورید. \\
	\textbf{پاسخ:}
	باید بررسی کنیم که ۱۱ را چگونه می‌توان به‌صورت جمع ضریبی از ۴ و ضریبی از ۳ نوشت:
	$$11=2 \times 4 + 3 \times 1 \Rightarrow A^{11} = (A^4)^2 \times (A^3)^1$$
	\[
	(A^4)^2=
	\begin{bmatrix}
		5 & 3 \\
		3 & 2
	\end{bmatrix}
	\begin{bmatrix}
		5 & 3 \\
		3 & 2
	\end{bmatrix}
	=
	\begin{bmatrix}
		34 & 21 \\
		21 & 13
	\end{bmatrix}
	\]
	\[
	A^{11}=(A^4)^2 \times A^3=
	\begin{bmatrix}
		34 & 21 \\
		21 & 13
	\end{bmatrix}
	\begin{bmatrix}
		3 & 2 \\
		2 & 1
	\end{bmatrix}
	=
	\begin{bmatrix}
		144 & 89 \\
		89 & 55
	\end{bmatrix}
	\]
\end{example}
\smallskip
\textbf{ماتریس خودتوان:}
ماتریس مربعی $A$ را خودتوان گویند هرگاه
$A^2=A$.
اگر یک ماتریس خودتوان باشد، آنگاه به ازای هر عدد طبیعی $n$ داریم
$A^n=A$.

\begin{example}
	تمام ماتریس‌هایی به صورت
	$A=
	\begin{bmatrix}
		a & 1 \\
		0 & b
	\end{bmatrix}$
	را بیابید که خودتوان باشند. \\
	\textbf{پاسخ:}
	\[
	A^2=
	\begin{bmatrix}
		a & 1 \\
		0 & b
	\end{bmatrix}
	\begin{bmatrix}
		a & 1 \\
		0 & b
	\end{bmatrix}
	=
	\begin{bmatrix}
		a^2 & a+b \\
		0 & b^2
	\end{bmatrix}
	\]
	\[
	A^2=A \Rightarrow
	\begin{bmatrix}
		a^2 & a+b \\
		0 & b^2
	\end{bmatrix}
	=
	\begin{bmatrix}
		a & 1 \\
		0 & b
	\end{bmatrix}
	\Rightarrow
	\begin{cases}
		a^2=a \\
		a+b=1 \\
		b^2=b
	\end{cases}
	\Rightarrow
	\begin{cases}
		a=0 \\
		b=1
	\end{cases}
	\text{یا} \
	\begin{cases}
		a=1 \\
		b=0
	\end{cases}
	\]
	بنابراین:
	\[
	A=
	\begin{bmatrix}
		0 & 1 \\
		0 & 1
	\end{bmatrix}
	\quad \text{یا} \quad
	A=
	\begin{bmatrix}
		1 & 1 \\
		0 & 0
	\end{bmatrix}
	\]
\end{example}

\begin{example}
	اگر $A$ خودتوان و $I$ ماتریس همانی هم‌مرتبه با $A$ باشد، حاصل
	$(A+I)^3$
	چیست؟ \\
	\textbf{پاسخ:}
	$A$ و $I$
	تعویض‌پذیرند. بنابراین:
	$$(A+I)^3 = A^3 + 3A^2 + 3A + I = A + 3A + 3A + I = 7A + I$$
\end{example}

\begin{theorem}
	اگر $A$ خودتوان و $I$ ماتریس همانی هم‌مرتبه با $A$ باشد، به ازای هر عدد طبیعی $n$ داریم:
	$$(A+I)^n = (2^n-1)A + I$$
\end{theorem}

\section{دترمینان}
به هر ماتریس مربعی می‌توان یک عدد نسبت داد که دترمینان آن ماتریس نامیده می‌شود. دترمینان یک ماتریس اطلاعات مفیدی راجع‌به خود ماتریس و خواص آن به ما خواهد داد.

اگر $A$ ماتریسی مربعی از مرتبه‌ی $n$ باشد، در این صورت دترمینان ماتریس $A$ را با نماد
$\det(A)=|A|$
نمایش می‌دهیم. برای ماتریس‌های مربعی مرتبه‌ی ۱ و ۲ داریم:
\[
A=[k]_{1 \times 1} \Rightarrow |A|=k
\]
\[
B=
\begin{bmatrix}
	a & b \\
	c & d \\	
\end{bmatrix}
\Rightarrow
|B|=ad-bc
\]

\begin{example}
	دترمینان ماتریس
	$A=\begin{bmatrix}
		\sin\alpha & \cos\alpha \\
		-\cos\alpha & \sin\alpha \\	
	\end{bmatrix}$
	را به دست آوردید.
	\[
	|A|=\sin^2\alpha + \cos^2\alpha = 1
	\]
\end{example}

پیش از تعریف دترمینان برای ماتریس مربعی مرتبه‌ی ۳ و بالاتر لازم است با چند تعریف دیگر آشنا شویم.
\\
\textbf{کهاد:}
فرض کنیم
$A=[a_{ij}]_{n \times n}$
ماتریسی دلخواه باشد. در این صورت $ij$-امین کهاد ماتریس $A$ را که با
$M_{ij}$
نمایش می‌دهیم، ماتریسی
$(n-1) \times (n-1)$
تعریف می‌کنیم که از حذف سطر $i$ و ستون $j$ ماتریس $A$ به دست می‌آید.

\begin{example}
	برای ماتریس
	$A=
	\begin{bmatrix}
		2 & -1 & 4 \\
		0 & 1 & 5 \\
		6 & 3 & -4
	\end{bmatrix}$،
	داریم
	$M_{13}=
	\begin{bmatrix}
		0 & 1 \\
		6 & 3
	\end{bmatrix}$
	و
	$M_{32}=
	\begin{bmatrix}
		2 & 4 \\
		0 & 5
	\end{bmatrix}$.
\end{example}
\bigskip
\textbf{همسازه:}
فرض کنیم
$A=[a_{ij}]_{n \times n}$
ماتریسی دلخواه باشد. در این صورت $ij$-امین همسازه‌ی ماتریس $A$ را که با
$C_{ij}$
نمایش می‌دهیم، به صورت زیر تعریف می‌کنیم:
$$C_{ij} = (-1)^{i+j} \left|M_{ij}\right|$$

\begin{example}
	ماتریس $A$ مثال قبل را در نظر بگیرید.
	\[
	C_{13} = (-1)^{1+3} \left|M_{13}\right| = (-1)^4
	\begin{vmatrix}
		0 & 1 \\
		6 & 3
	\end{vmatrix}
	= 1 \times (0-6) = -6
	\]
	\[
	C_{32} = (-1)^{3+2} \left|M_{32}\right| = (-1)^5
	\begin{vmatrix}
		2 & 4 \\
		0 & 5
	\end{vmatrix}
	= (-1) \times (10 - 0) = -10
	\]
\end{example}

\begin{theorem} \label{thmdet3}
	فرض کنیم
	$A=[a_{ij}]_{n \times n}$
	ماتریسی دلخواه باشد. در این صورت حاصل عددی فرمول‌های
	\LR{\begin{itemize}
			\item $a_{11}C_{11} + a_{12}C_{12} + \cdots + a_{1n}C_{1n}$
			\item $a_{21}C_{21} + a_{22}C_{22} + \cdots + a_{2n}C_{2n}$
			\\ \vdots
			\item $a_{n1}C_{n1} + a_{n2}C_{n2} + \cdots + a_{nn}C_{nn}$
			\item $a_{11}C_{11} + a_{21}C_{21} + \cdots + a_{n1}C_{n1}$
			\item $a_{12}C_{12} + a_{22}C_{22} + \cdots + a_{n2}C_{n2}$
			\\ \vdots
			\item $a_{1n}C_{1n} + a_{2n}C_{2n} + \cdots + a_{nn}C_{nn}$
	\end{itemize}}
	همگی با هم مساوی هستند.
\end{theorem}

\textbf{تعریف دترمینان:}
فرض کنیم $A$ ماتریسی مربعی و دلخواه باشد.
$|A|$
را یکی از اعداد مساوی معرفی‌شده در قضیه‌ی
\ref{thmdet3}
تعریف می‌کنیم. برای محاسبه‌ی دترمینان کافی است از یکی از این فرمول‌ها استفاده کنیم.

\begin{example}
	فرض کنید اعداد حقیقی $a$، $b$ و $c$ بزرگتر از صفر و مخالف یکدیگر باشند. داریم:
	\[
	\det(A)=
	\begin{vmatrix}
		a & b & c \\
		b & c & a \\
		c & a & b
	\end{vmatrix} = 0
	\]
	ثابت کنید که
	$a^2+b^2+c^2 = ab+ac+bc$.
	\\ \textbf{پاسخ:}
	\[
	C_{11}=(-1)^{1+1}
	\begin{vmatrix}
		c & a \\
		a & b
	\end{vmatrix}
	\qquad
	C_{12}=(-1)^{1+2}
	\begin{vmatrix}
		b & a \\
		c & b
	\end{vmatrix}
	\qquad
	C_{13}=(-1)^{1+3}
	\begin{vmatrix}
		b & c \\
		c & a
	\end{vmatrix}
	\]
	\begin{align*}
		|A| &= a_{11}C_{11} + a_{12}C_{12} + a_{13}C_{13} \\
		&= a(bc-a^2) -b(b^2-ac) +c(ab-c^2) \\
		&= 3abc-a^3+b^3+c^3 \\
		&= -(a+b+c)(a^2+b^2+c^2-ab-ac-bc) = 0
	\end{align*}
	چون
	$a+b+c \neq 0$،
	بنابراین:
	$$a^2+b^2+c^2-ab-ac-bc = 0 \Rightarrow a^2+b^2+c^2 = ab+ac+bc$$
\end{example}

\begin{example}
	مطلوب است محاسبه‌ی دترمینان
	\[
	A=
	\begin{bmatrix}
		1 & 1 & 1 & 1 \\
		0 & 1 & 2 & 3 \\
		1 & 2 & 3 & 4 \\
		2 & 3 & 4 & 5
	\end{bmatrix}
	\]
	\textbf{پاسخ:}
	\[
	\det(A)= a_{11}C_{11} + a_{12}C_{12} + a_{13}C_{13} + a_{14}C_{14}
	\]
	\[
		\det(A) =
		1 \times
		\begin{vmatrix}
			1 & 2 & 3 \\
			2 & 3 & 4 \\
			3 & 4 & 5
		\end{vmatrix}
		-1 \times
		\begin{vmatrix}
			0 & 2 & 3 \\
			1 & 3 & 4 \\
			2 & 4 & 5
		\end{vmatrix}
		+1 \times
		\begin{vmatrix}
			0 & 1 & 3 \\
			1 & 2 & 4 \\
			2 & 3 & 5
		\end{vmatrix}
		-1 \times
		\begin{vmatrix}
			0 & 1 & 2 \\
			1 & 2 & 3 \\
			2 & 3 & 4
		\end{vmatrix}
		= 0
	\]
\end{example}

\begin{example}
	اگر $A$ ماتریسی
	$3 \times 3$
	باشد و
	$|A|=5$،
	در این صورت حاصل
	$\left||A|A\right|$
	را بیابید. \\
	\textbf{پاسخ:}
\end{example}

\begin{example}
	فرض کنید $A$ و $B$ دو ماتریس
	$3 \times 3$
	باشند که $A$ متقارن است. ثابت کنید
	$|A+B|=|A+B^t|$.
	\\ \textbf{پاسخ:}
\end{example}

\subsection{ویژگی‌های دترمینان}
در این قسمت، ویژگی‌های دترمینان را بیان می‌کنیم. برای سادگی، ویژگی‌ها را برای ماتریس مربعی مرتبه‌ی ۲ ذکر می‌کنیم.
\begin{enumerate}
	\item 
	اگر همه‌ی درایه‌های یک سطر (ستون) را در عدد حقیقی $r$ ضرب کنیم، آنگاه دترمینان ماتریس جدید $r$ برابر ماتریس قبلی است.
	\[
	\begin{vmatrix}
		ra & b \\
		rc & d
	\end{vmatrix}
	= r
	\begin{vmatrix}
		a & b \\
		c & d
	\end{vmatrix}
	\]
	\item 
	دترمینان ماتریس قطری برابر با حاصل‌ضرب درایه‌های قطر اصلی آن است.
	\[
	\begin{vmatrix}
		a & 0 \\
		0 & b
	\end{vmatrix} = ab
	\]
	\item 
	اگر همه‌ی درایه‌های یک سطر (ستون) برابر صفر باشند، آنگاه دترمینان ماتریس برابر صفر است.
	\[
	\begin{vmatrix}
		0 & 0 \\
		c & d
	\end{vmatrix} = 0
	\]
	\item 
	اگر دو سطر (دو ستون) یکسان باشد، آنگاه دترمینان ماتریس برابر صفر است.
	\[
	\begin{vmatrix}
		a & b \\
		a & b
	\end{vmatrix} = 0
	\]
	\item 
	اگر یک سطر (ستون) مضربی از یک سطر (ستون) دیگر باشد، آنگاه دترمینان ماتریس برابر صفر است.
	\[
	\begin{vmatrix}
		a & b \\
		ka & kb
	\end{vmatrix} = 0
	\]
	\item 
	اگر مضرب‌های یک سطر (ستون) را به سطر (ستون) دیگر اضافه کنیم، مقدار دترمینان تغییر نمی‌کند.
	\[
	\begin{vmatrix}
		a & b \\
		c & d
	\end{vmatrix}
	=
	\begin{vmatrix}
		a & b+ma \\
		c & d+mc
	\end{vmatrix}
	\]
	\item 
	اگر جای دو سطر (دو ستون) را عوض کنیم، آنگاه دترمینان ماتریس جدید قرینه‌ی دترمینان ماتریس قبلی است.
	\[
	\begin{vmatrix}
		c & d \\
		a & b
	\end{vmatrix}
	= -
	\begin{vmatrix}
		a & b \\
		c & d
	\end{vmatrix}
	\]
	\item 
	دترمینان ترانهاده‌ی یک ماتریس با دترمینان خود آن ماتریس برابر است.
	\[
	\begin{vmatrix}
		a & b \\
		c & d
	\end{vmatrix}
	=
	\begin{vmatrix}
		a & c \\
		b & d
	\end{vmatrix}
	\]
	\item 
	اگر همه‌ی درایه‌های یک سطر (ستون)، مجموع دو مقدار باشند، می‌توان دترمینان را برابر مجموع دو دترمینان نوشت.
	\[
	\begin{vmatrix}
		a+b & c+d \\
		e & f
	\end{vmatrix}
	=
	\begin{vmatrix}
		a & c \\
		e & f
	\end{vmatrix}
	+
	\begin{vmatrix}
		b & d \\
		e & f
	\end{vmatrix}
	\]
\end{enumerate}

\begin{theorem}
	اگر
	$A=[a_{ij}]_{m \times m}$ و
	$B=[b_{ij}]_{m \times m}$،
	آنگاه
	$|AB|=|A||B|$. \\
	\textbf{نتیجه:}
	فرض کنیم
	$A=[a_{ij}]_{m \times m}$
	و
	$n \in \mathbb{N}$.
	در این صورت
	$|A^n|=|A|^n$.
\end{theorem}

\begin{theorem}
	اگر
	$A=[a_{ij}]_{n \times n}$
	و
	$r \in \mathbb{R}$،
	در این صورت
	$|rA| = r^n|A|$.
\end{theorem}

\begin{example}
	به کمک ویژگی‌های دترمینان‌ها ثابت کنید
	\[
	\begin{vmatrix}
		1 & 1 & 1 \\
		x & y & z \\
		x^2 & y^2 & z^2
	\end{vmatrix}
	= (y-x)(z-x)(z-y)
	\]
	\\ \textbf{پاسخ:}
\end{example}

\begin{example}
	به کمک ویژگی‌های دترمینان‌ها ثابت کنید
	\[
	\begin{vmatrix}
		1+x & y & z \\
		x & 1+y & z \\
		x & y & 1+z
	\end{vmatrix}
	= 1+x+y+z
	\]
	\\ \textbf{پاسخ:}
\end{example}

\section{ماتریس وارون}
فرض کنیم
$A=[a_{ij}]_{n \times n}$.
وارون ماتریس $A$ (در صورت وجود) ماتریسی است چون
$B=[b_{ij}]_{n \times n}$
به طوری که
$AB=BA=I_n$.
در این صورت $B$ را وارون $A$ می‌نامیم و با
$A^{-1}$
نشان می‌دهیم.

\begin{theorem}
	فرض کنیم $A$ ماتریسی مربعی باشد که وارون‌پذیر است. در این صورت وارون $A$ منحصربه‌فرد است.
\end{theorem}

\begin{theorem}
	فرض کنیم $A$ ماتریسی مربعی باشد که وارون‌پذیر است. در این صورت
	$|A| \neq 0$.
\end{theorem}

\begin{example}
	ماتریس
	$A=
	\begin{bmatrix}
		2 & -1 \\
		3 & -4
	\end{bmatrix}$،
	$I$
	ماتریس همانی هم‌مرتبه با $A$ و
	$\alpha$ و $\beta$
	دو عدد حقیقی هستند، به‌طوری که
	$\alpha A^{-1} + \beta I = A$.
	مقدار
	$\alpha$ و $\beta$
	را به دست آورید. \\
	\textbf{پاسخ:}
	دو طرف تساوی را در ماتریس $A$ ضرب می‌کنیم.
	$$\alpha A^{-1}A + \beta IA = AA \Rightarrow \alpha I + \beta A = A^2$$
	ماتریس
	$A^2$
	را حساب می‌کنیم و در تساوی بالا جایگذاری می‌کنیم.
	\[
	A^2 = 
	\begin{bmatrix}
		2 & -1 \\
		3 & -4
	\end{bmatrix}
	\begin{bmatrix}
		2 & -1 \\
		3 & -4
	\end{bmatrix}
	=
	\begin{bmatrix}
		1 & 2 \\
		-6 & 13
	\end{bmatrix}
	\]
	\[
	\alpha
	\begin{bmatrix}
		1 & 0 \\
		0 & 1
	\end{bmatrix}
	+ \beta
	\begin{bmatrix}
		2 & -1 \\
		3 & -4
	\end{bmatrix}
	=
	\begin{bmatrix}
		1 & 2 \\
		-6 & 13
	\end{bmatrix}
	\Rightarrow
	\begin{bmatrix}
		\alpha + 2\beta & -\beta \\
		3\beta & \alpha -4\beta
	\end{bmatrix}
	=
	\begin{bmatrix}
		1 & 2 \\
		-6 & 13
	\end{bmatrix}
	\]
	\[
	\Rightarrow
	\begin{cases}
		\alpha + 2\beta = 1 \\
		-\beta = 2
	\end{cases}
	\Rightarrow
	\alpha = 5 \ , \ \beta = -2
	\]
\end{example}

\subsection{وارون ماتریس مربعی مرتبه‌ی ۲}

\begin{theorem}
	ماتریس
	$A=
	\begin{bmatrix}
		a_{11} & a_{12} \\
		a_{21} & a_{22}
	\end{bmatrix}$
	وارون‌پذیر است اگر و تنها اگر
	$|A| \neq 0$.
	در این حالت داریم:
	\[
	A^{-1} = \frac{1}{|A|}
	\begin{bmatrix}
		a_{22} & -a_{12} \\
		-a_{21} & a_{11}
	\end{bmatrix}
	\]
\end{theorem}

\begin{example}
	وارون ماتریس
	$A=
	\begin{bmatrix}
		10 & 2 \\
		4 & 1
	\end{bmatrix}$
	به دست آورید. \\
	\textbf{پاسخ:}
	چون
	$|A|=2 \neq 0$
	پس $A$ وارون‌پذیر است و داریم:
	\[
	A^{-1}=\frac{1}{2}
	\begin{bmatrix}
		1 & -2 \\
		-4 & 10
	\end{bmatrix}
	=
	\begin{bmatrix}
		\dfrac{1}{2} & -1 \\
		-2 & 5
	\end{bmatrix}
	\]
\end{example}

\begin{example}
	اگر
	$A=
	\begin{bmatrix}
		4 & 3 \\
		2 & 5
	\end{bmatrix}$
	و
	$B=
	\begin{bmatrix}
		-2 & -3 \\
		5 & -1
	\end{bmatrix}$،
	حاصل عبارت
	$(2A^{-1}-3B^{-1})$
	را بیابید. \\
	\textbf{پاسخ:}
\end{example}

\begin{example}
	اگر
	$A=
	\begin{bmatrix}
		5 & 2 \\
		3 & 2
	\end{bmatrix}$،
	ابتدا ماتریس
	$A^{-1}$
	را به دست آورده و $|A|$ را با $|A^{-1}|$ مقایسه کنید. \\
	\textbf{پاسخ:}
\end{example}

\subsection{وارون ماتریس مربعی مرتبه‌ی ۳}

\begin{theorem}
	ماتریس
	$A=
	\begin{bmatrix}
		a_{11} & a_{12} & a_{13} \\
		a_{21} & a_{22} & a_{23} \\
		a_{31} & a_{32} & a_{33}
	\end{bmatrix}$
	وارون‌پذیر است اگر و تنها اگر
	$|A| \neq 0$.
	در این حالت داریم:
	\[
	A^{-1} = \frac{1}{|A|}
	\begin{bmatrix}
		C_{11} & C_{21} & C_{31} \\
		C_{12} & C_{22} & C_{32} \\
		C_{13} & C_{23} & C_{33}
	\end{bmatrix}
	\]
\end{theorem}

\begin{example}
	برای ماتریس
	$A=
	\begin{bmatrix}
		2 & 3 & -4 \\
		0 & -4 & 2 \\
		1 & -1 & 5
	\end{bmatrix}$
	داریم:
	\[
	C_{11}=
	\begin{vmatrix}
		-4 & 2 \\
		-1 & 5
	\end{vmatrix} = -18
	\ , \quad
	C_{12}= -
	\begin{vmatrix}
		0 & 2 \\
		1 & 5
	\end{vmatrix} = 2
	\ , \quad
	C_{13}=
	\begin{vmatrix}
		0 & -4 \\
		1 & -1
	\end{vmatrix} = 4
	\ ,
	\]
	\[
	C_{21}= -
	\begin{vmatrix}
		3 & -4 \\
		-1 & 5
	\end{vmatrix} = -11
	\ , \quad
	C_{22}=
	\begin{vmatrix}
		2 & -4 \\
		1 & 5
	\end{vmatrix} = 14
	\ , \quad
	C_{23}= -
	\begin{vmatrix}
		2 & 3 \\
		1 & -1
	\end{vmatrix} = 5
	\ ,
	\]
	\[
	C_{31}=
	\begin{vmatrix}
		3 & -4 \\
		-4 & 2
	\end{vmatrix} = -10
	\ , \quad
	C_{32}= -
	\begin{vmatrix}
		2 & -4 \\
		0 & 2
	\end{vmatrix} = -4
	\ , \quad
	C_{33}=
	\begin{vmatrix}
		2 & 3 \\
		0 & -4
	\end{vmatrix} = -8
	\ ,
	\]
	همچنین
	$|A|=-46$،
	بنابراین:
	\[
	A^{-1} = -\frac{1}{46}
	\begin{bmatrix}
		-18 & -11 & -10 \\
		2 & 14 & -4 \\
		4 & 5 & -8
	\end{bmatrix}
	=
	\begin{bmatrix}
		\dfrac{9}{23} & \dfrac{11}{46} & \dfrac{5}{23} \\
		\dfrac{-1}{23} & \dfrac{-7}{23} & \dfrac{2}{23} \\
		\dfrac{-2}{23} & \dfrac{-5}{46} & \dfrac{4}{23}
	\end{bmatrix}
	\]
\end{example}

\begin{example}[پیش‌نیاز: دستگاه معادلات خطی]
	دو ریاضی‌دان به کمک یک ماتریس رمزنگاری به هم پیام می‌دهند. می‌دانید ماتریس رمزنگاری آن‌ها یک ماتریس ۳ در ۳ است. به هر حرف الفبای انگلیسی به ترتیب اعداد ۱ تا ۲۶ و به فاصله‌ی بین دو کلمه عدد ۲۷ اختصاص داده می‌شود. شما سه پیام لورفته‌ی آن‌ها را در اختیار دارید:
	\[
	HEY^*=
	\begin{bmatrix}
		11 \\ 29 \\ -6
	\end{bmatrix} \ ,
	\quad
	TWO^*=
	\begin{bmatrix}
		117 \\ 83 \\ -94
	\end{bmatrix} \ ,
	\quad
	MAX^*=
	\begin{bmatrix}
		6 \\ 40 \\ -5
	\end{bmatrix}
	\]
	اکنون پیامی به صورت زیر مخابره شده است. آن را رمزگشایی کنید.
	\[
	\begin{bmatrix}
		62 & 99 & 51 \\
		39 & 48 & 17 \\
		-47 & -72 & -37
	\end{bmatrix}
	\]
	\\
	\textbf{پاسخ:}
	ابتدا حروف سه پیام اول را مطابق الفبا به عدد تبدیل می‌کنیم.
	\[
	HEY=
	\begin{bmatrix}
		8 \\ 5 \\ 25
	\end{bmatrix} \ ,
	\quad
	TWO=
	\begin{bmatrix}
		20 \\ 23 \\ 15
	\end{bmatrix} \ ,
	\quad
	MAX=
	\begin{bmatrix}
		13 \\ 1 \\ 24
	\end{bmatrix}
	\]
	ماتریس رمزنگاری را به صورت زیر فرض می‌کنیم:
	\[
	K=
	\begin{bmatrix}
		a & b & c \\
		d & e & f \\
		g & h & i
	\end{bmatrix}
	\]
	داریم:
	\[
	\begin{bmatrix}
		a & b & c \\
		d & e & f \\
		g & h & i
	\end{bmatrix}
	\begin{bmatrix}
		8 \\ 5 \\ 25
	\end{bmatrix}
	=
	\begin{bmatrix}
		11 \\ 29 \\ -6
	\end{bmatrix}
	\Rightarrow
	\begin{bmatrix}
		8a+5b+25c \\
		8d+5e+25f \\
		8g+5h+25i
	\end{bmatrix}
	=
	\begin{bmatrix}
		11 \\ 29 \\ -6
	\end{bmatrix}
	\]
	
	\[
	\begin{bmatrix}
		a & b & c \\
		d & e & f \\
		g & h & i
	\end{bmatrix}
	\begin{bmatrix}
		20 \\ 23 \\ 15
	\end{bmatrix}
	=
	\begin{bmatrix}
		117 \\ 83 \\ -94
	\end{bmatrix}
	\Rightarrow
	\begin{bmatrix}
		20a+23b+15c \\
		20d+23e+15f \\
		20g+23h+15i
	\end{bmatrix}
	=
	\begin{bmatrix}
		117 \\ 83 \\ -94
	\end{bmatrix}
	\]
	
	\[
	\begin{bmatrix}
		a & b & c \\
		d & e & f \\
		g & h & i
	\end{bmatrix}
	\begin{bmatrix}
		13 \\ 1 \\ 24
	\end{bmatrix}
	=
	\begin{bmatrix}
		6 \\ 40 \\ -5
	\end{bmatrix}
	\Rightarrow
	\begin{bmatrix}
		13a+b+24c \\
		13d+e+24f \\
		13g+h+24i
	\end{bmatrix}
	=
	\begin{bmatrix}
		6 \\ 40 \\ -5
	\end{bmatrix}
	\]
	با حل دستگاه بالا خواهیم داشت:
	$$a=2,\ b=4,\ c=-1, \ d=3, \ e=1, \ f=0, \ g=-2, \ h=-3, \ i=1$$
	\[
	K=
	\begin{bmatrix}
		2 & 4 & -1 \\
		3 & 1 & 0 \\
		-2 & -3 & 1
	\end{bmatrix}
	\]
\end{example}

\subsection{خواص ماتریس وارون}
فرض کنیم $A$ و $B$ وارون‌پذیر و $c$ یک عدد غیرصفر و $n$ یک عدد طبیعی باشد. آنگاه:
\LR{
	\begin{enumerate}
		\item $(A^{-1})^{-1}=A$
		\item $(cA)^{-1}=\frac{1}{c}A^{-1}$
		\item $(AB)^{-1}=B^{-1}A^{-1}$
		\item $(A^n)^{-1}=(A^{-1})^n$
		\item $(A^t)^{-1}=(A^{-1})^t$
\end{enumerate}}
\vspace*{-0.5cm}

\begin{example}
	اگر
	$A^{-1}=
	\begin{bmatrix}
		1 & -1 \\
		-3 & 4
	\end{bmatrix}$،
	آنگاه
	$(A^t)^{-1}$
	را به دست آورید. \\
	\textbf{پاسخ:}
	\[
	(A^t)^{-1}=(A^{-1})^t=
	\begin{bmatrix}
		1 & -1 \\
		-3 & 4
	\end{bmatrix}^t
	=
	\begin{bmatrix}
		1 & -3 \\
		-1 & 4
	\end{bmatrix}
	\]
	
\end{example}

\section{مسائل}

\begin{problem} \label{prmat1}
	$A$
	یک ماتریس پادمتقارن مرتبه‌ی ۳ با درایه‌های صحیح است، به طوری که در دو شرط زیر صدق می‌کند: \\
	۱) حاصل‌ضرب همه‌ی درایه‌های غیر قطر اصلی آن $-4$ است. \\
	۲) حاصل‌جمع درایه‌های ستون سوم آن ۳ است. \\
	$A$
	را مشخص کنید.
\end{problem}

\begin{problem} \label{prmat2}
	$AB$ و $AC$
	را به دست آورید و با هم مقایسه کنید.
	\[
	A=
	\begin{bmatrix}
		1 & 2 & 0 \\
		1 & 1 & 0 \\
		-1 & 4 & 0
	\end{bmatrix},
	\qquad
	B=
	\begin{bmatrix}
		1 & 2 & 3 \\
		1 & 1 & -1 \\
		11 & 10 & 1
	\end{bmatrix},
	\qquad
	C=
	\begin{bmatrix}
		1 & 2 & 3 \\
		1 & 1 & -1 \\
		12 & 2 & -4
	\end{bmatrix}
	\]
\end{problem}

\begin{problem} \label{prmat3}
	ماتریس‌های زیر را با درایه‌هایشان مشخص کنید. سپس حاصل $AB$ را به دست آورید.
	\[
	A = [a_{ij}]_{4 \times 3} \quad , \quad a_{ij} =
	\begin{cases}
		i+j \qquad i=j \\
		i-j \qquad i>j \\
		j-i \qquad j>i
	\end{cases}
	\]
	\[
	B = [b_{ij}]_{3 \times 2} \quad , \quad b_{ij} =
	\begin{cases}
		i^j \qquad i=j \\
		j^i \qquad i \neq j
	\end{cases}
	\]
\end{problem}

\begin{problem} \label{prmat4}
	با فرض
	$A=
	\begin{bmatrix}
		2 & 1 & 1 \\
		1 & 2 & 1 \\
		1 & 1 & 2
	\end{bmatrix}$،
	حاصل
	$A^2 - 4A - 5I$
	را بیابید.
\end{problem}

\begin{problem} \label{prmat5}
	$A$ و $B$
	دو ماتریس مربعی هم‌مرتبه‌اند و عمل $\circledast$ به صورت زیر تعریف شده است:
	$$A \circledast B = AB - BA$$
	نشان دهید برای ماتریس مربعی $C$ که با $A$ و $B$ هم‌مرتبه باشد، داریم:
	$$(A \circledast B) \circledast C + (B \circledast C) \circledast A + (C \circledast A) \circledast B = O$$
\end{problem}

\begin{problem} \label{prmat6}
	اگر
	$A=
	\begin{bmatrix}
		30 & 25 \\
		-36 & -30
	\end{bmatrix}$،
	ماتریس
	$A^{100}$
	را به دست آورید.
\end{problem}

\begin{problem} \label{prmat7}
	با چه شرطی، این برابری برقرار است؟
	\[
	\begin{vmatrix}
		1 & \cos\alpha & \cos\beta \\
		\cos\alpha & 1 & \cos\gamma \\
		\cos\beta & \cos\gamma & 1
	\end{vmatrix}
	=
	\begin{vmatrix}
		0 & \cos\alpha & \cos\beta \\
		\cos\alpha & 0 & \cos\gamma \\
		\cos\beta & \cos\gamma & 0
	\end{vmatrix}
	\]
\end{problem}

\begin{problem} \label{prmat8}
	به کمک ویژگی‌های دترمینان‌ها ثابت کنید:
	\[
	\begin{vmatrix}
		1 & 1+x & x^2(y+z) \\
		1 & 1+y & y^2(x+z) \\
		1 & 1+z & z^2(x+y)
	\end{vmatrix} = 0
	\]
\end{problem}

\begin{problem} \label{prmat9}
	به کمک ویژگی‌های دترمینان‌ها ثابت کنید:
	\[
	\begin{vmatrix}
		yz & x^2 & x^2 \\
		y^2 & xz & y^2 \\
		z^2 & z^2 & xy
	\end{vmatrix}
	=
	\begin{vmatrix}
		yz & xy & xz \\
		xy & xz & yz \\
		xz & yz & xy
	\end{vmatrix}
	\]
\end{problem}

\begin{problem}[پیش‌نیاز: احتمال] \label{prmat10}
	ماتریس‌های
	$A=
	\begin{bmatrix}
		a & -7 \\
		-2 & 0
	\end{bmatrix}$
	و
	$B=
	\begin{bmatrix}
		0 & b \\
		1 & 1
	\end{bmatrix}$
	مفروض‌اند. یک تاس را دو بار پرتاب کرده و به جای $a$ و $b$ در ماتریس‌های $A$ و $B$ جای‌گذاری می‌کنیم. چقدر احتمال دارد مجموع دو ماتریس وارون‌پذیر باشد؟
\end{problem}

\begin{problem} \label{prmat11}
	ماتریس مربعی $A$ پوچ‌توان نامیده می‌شود، در صورتی که به ازای عددی طبیعی مانند $n$ داشته باشیم
	$A^n=O$.
	کوچک‌ترین عدد طبیعی واجد این خاصیت را اندیس $A$ (یا مرتبه‌ی پوچی $n$) گویند. نشان دهید ماتریس زیر پوچ‌توان با اندیس ۳ است.
	\[
	A=
	\begin{bmatrix}
		1 & 1 & 3 \\
		5 & 2 & 6 \\
		-2 & -1 & -3
	\end{bmatrix}
	\]
\end{problem}

\begin{problem} \label{prmat12}
	اگر $A$ ماتریسی وارون‌پذیر باشد، به‌طوری که
	$\sqrt{3}A =
	\begin{bmatrix}
		-2 & |A|+1 \\
		3|A| & -|A|
	\end{bmatrix}$،
	آنگاه $|A|$ را به دست آورید.
\end{problem}
